\documentclass[10pt,letterpaper]{article}

\usepackage[T1]{fontenc}
\usepackage[utf8]{inputenc}
\usepackage{lmodern}
\usepackage[margin=0.82in,headheight=14pt]{geometry}
\usepackage{microtype,amsmath,amssymb,array,tabularx,longtable,booktabs,multicol}
\usepackage{graphicx,xcolor,enumitem,fancyhdr,hyperref,titlesec,needspace}
\definecolor{Navy}{HTML}{122A42}
\definecolor{Ocean}{HTML}{147D91}
\definecolor{Soft}{HTML}{E9F2F5}
\definecolor{Ink}{HTML}{263745}
\definecolor{Gray}{HTML}{5F6873}
\hypersetup{colorlinks=true,linkcolor=Ocean,urlcolor=Ocean,bookmarksnumbered=true}
\setlength{\parindent}{0pt}
\emergencystretch=2em
\setlength{\parskip}{5.3pt}
\setlist{topsep=3pt,itemsep=2pt,parsep=0pt,leftmargin=*}
\titleformat{\section}{\large\bfseries\color{Navy}}{\thesection}{0.65em}{}
\titleformat{\subsection}{\normalsize\bfseries\color{Ocean}}{\thesubsection}{0.6em}{}
\titleformat{\subsubsection}{\normalsize\bfseries\color{Navy}}{\thesubsubsection}{0.5em}{}
\pagestyle{fancy}\fancyhf{}
\fancyhead[L]{\textsf{Architecture Documentation \& Engineering}}
\fancyhead[R]{\textsf{\DocCategory}}
\fancyfoot[L]{\footnotesize\textcolor{Gray}{Controlled learning / engineering artifact}}
\fancyfoot[R]{\footnotesize\thepage}
\renewcommand{\headrulewidth}{0.35pt}
\newcommand{\DocCategory}{Suite}
\newcommand{\ArchTitle}[4]{%
\renewcommand{\DocCategory}{#3}%
{\small\sffamily\bfseries\color{Ocean} #3 \hfill #4}\par\vspace{7pt}
{\Large\bfseries\color{Navy}#1}\par\vspace{4pt}
{\normalsize\color{Gray}#2}\par\vspace{11pt}
\textcolor{Ocean}{\rule{\linewidth}{1.15pt}}\vspace{4pt}
}
\newcommand{\Note}[1]{\par\smallskip\noindent\fcolorbox{Ocean}{Soft}{\parbox{0.94\linewidth}{\small\textbf{Engineering note.}~#1}}\par\smallskip}
\newcommand{\Outcome}[1]{\textbf{Outcome.} #1\par}
\newcommand{\Evidence}[1]{\textbf{Required evidence.} #1\par}
\newcommand{\BlankLines}[1]{\par\begingroup\setlength{\parskip}{9pt}\count255=0\loop\ifnum\count255<#1\noindent\rule{\linewidth}{0.35pt}\par\advance\count255 by 1\repeat\endgroup\vspace{4pt}}
\newcommand{\BlankBox}[1]{\noindent\fbox{\parbox[t][#1][t]{0.98\linewidth}{\mbox{}}}\par\smallskip}
\newcommand{\Field}[2]{\Needspace{3\baselineskip}\textbf{#1}\par\textit{\color{Gray}#2}\par\BlankLines{2}}
\newcommand{\FieldBlank}[1]{\Needspace{3\baselineskip}\textbf{#1}\par\BlankLines{2}}
\newcommand{\DocProvenance}{\textit{Source foundation: Clements et al., \emph{Documenting Software Architectures: Views and Beyond}, 2nd ed. (2011). Interpretive examples, exercises, checklists, and templates are original adaptations, not reproductions.}\par}
\newcommand{\SmallHeader}[1]{\smallskip\textbf{\color{Navy}#1}\par}
\newcolumntype{Y}{>{\raggedright\arraybackslash}X}
\newcommand{\Tick}{\(\square\)\,}

\begin{document}
\ArchTitle{Study Guide 10: Building the Documentation Package}{Reading, architecture reasoning, implementation exercise, and evaluation}{Study Guide}{ARCH-SG-10}
\DocProvenance
\section{Conceptual focus}
A view comprises presentation plus element catalog, context, variability, and rationale; the package adds cross-view information.

\section{Reading map}
Book location: Building the Documentation Package. Principal topics: standard view template, documentation beyond views, correspondence, packaging. Read the chapter discussion, then its summary checklist and selected discussion questions. Record terms precisely before adapting them to an implementation.

\section{Learning objectives}
\begin{itemize}
\item Explain standard view template, documentation beyond views, correspondence, packaging in terms of a stakeholder question.
\item Create reader roadmap and mapping between views; control revisions and glossary.
\item Separate book-supported architectural concepts from optional implementation techniques such as Terraform, OpenAPI, and CI-generated evidence.
\end{itemize}
\section{Engineering interpretation}
Create reader roadmap and mapping between views; control revisions and glossary.

\section{Worked laboratory assignment}
Assemble a three-view architecture description with consistent IDs and map to requirements.

\section{Evidence to submit}
Complete draft AD with mappings

\section{Review prompts}
\begin{itemize}
\item What information belongs outside any single view?
\item What relation types or behavior are intentionally outside the scope of this view?
\item Which system change would invalidate your claims, and how would you detect it?
\end{itemize}
\section{Self-assessment rubric}
\small
\begin{tabularx}{\linewidth}{@{}YYYY@{}}
\toprule
\textbf{Criterion} & \textbf{Beginning (0)} & \textbf{Working (1)} & \textbf{Proficient (2)} \\
\midrule
Concepts & Undefined or conflated & Mostly distinguished & Accurate, precise and consistently applied \\
Architecture evidence & Diagram only & Evidence incomplete & Typed model with catalog and provenance \\
Stakeholder fitness & No stated reader & Reader stated & Reader can answer target question \\
\bottomrule
\end{tabularx}
\normalsize
\section{Technical concept-to-evidence crosswalk}
\small
\begin{tabularx}{\linewidth}{@{}YYY@{}}
\toprule
\textbf{Concept} & \textbf{Engineering interpretation} & \textbf{Verification evidence} \\
\midrule
Primary presentation & The first view page selects important elements and typed relations; supporting material explains it. & Run a legend interpretability check. \\
Element catalog & Properties, interfaces and behavior fill gaps left by primary diagrams. & Join every drawn ID to a catalog entry. \\
Beyond views & Roadmap, overview, mapping, rationale and directory improve the document as a whole. & Trace one reader question across two view families. \\
Packaging and revision & Versioned releases and controlled links prevent accidentally mixing contradictory artifacts. & Define a baseline and release readiness checklist. \\
\bottomrule
\end{tabularx}
\normalsize
\section{Chapter-specific analysis tasks}
\begin{itemize}
\item Identify and define each relevant element and relation before you choose a diagram notation.
\item Produce one counterexample showing how a plausible but incorrect model could mislead a stakeholder.
\item Identify the source and timestamp of at least one fact about the implementation; label inference separately.
\end{itemize}
\section{Links to suite artifacts}
Use the appropriate architecture view, viewpoint, quality scenario, traceability workbook, and review templates. Preserve stable IDs across artifacts. Record your chapter learning reflections in the curriculum log.


\end{document}
